% ----------------------------------------------------------------------------
% FemPlate — technical documentation
%
% Build:  pdflatex technical.tex  (twice, for references)
%
% (c) 2026 Olivier Doaré — LGPL-3.0-or-later
% ----------------------------------------------------------------------------
\documentclass[11pt,a4paper]{article}

\usepackage[T1]{fontenc}
\usepackage[utf8]{inputenc}
\usepackage{amsmath,amssymb,bm}
\usepackage[margin=2.6cm]{geometry}
\usepackage{booktabs}
\usepackage{hyperref}
\usepackage{xcolor}
\hypersetup{colorlinks=true, linkcolor=blue!50!black, citecolor=blue!50!black,
            urlcolor=blue!50!black}

\newcommand{\lap}{\Delta}
\newcommand{\bilap}{\Delta^2}
\newcommand{\dd}{\mathrm{d}}
\newcommand{\code}[1]{\texttt{#1}}

\title{FemPlate: finite-element modal synthesis of plates and membranes
       of arbitrary shape\\[1ex]
       \large Technical documentation}
\author{Olivier Doaré\thanks{FX-Mechanics, \url{github.com/odoare}. Model,
        implementation and documentation developed with Claude (Anthropic).}}
\date{July 2026 --- FemPlate v0.1}

\begin{document}
\maketitle

\begin{abstract}
FemPlate is an audio plugin that simulates the linear and weakly nonlinear
vibration of a thin plate (or tensioned membrane) of arbitrary user-drawn
shape, with mixed boundary conditions. The plate eigenmodes are computed by
a non-conforming (Morley) finite-element discretisation and a dense subspace
eigensolver, both implemented in the reusable, JUCE-free library
\code{FxmeTools/acoustics}. Sound is produced by a bank of resonant
band-pass filters, one per mode, following the modal-synthesis approach of
the MechanOdd project. A first-order tension law allows audio-rate retuning
of the whole bank without re-solving the eigenproblem, and this same
mechanism supports an efficient approximation of geometric (von
K\'arm\'an--Berger) nonlinearity: amplitude-dependent pitch glide and a
cubic inter-mode energy-cascade surrogate, both realtime-safe and
unconditionally stable. This document details the mathematical model, the
discretisation, the eigensolution, the audio implementation and the
nonlinear approximations, with references to the literature.
\end{abstract}

\tableofcontents

% ============================================================================
\section{Continuous model}
\label{sec:model}

\subsection{Scaled Kirchhoff plate with tension}

The transverse deflection $w(\bm{x},t)$ of a thin plate occupying a domain
$\Omega\subset\mathbb{R}^2$ is modelled by Kirchhoff--Love theory
\cite{leissa,timoshenko} augmented by a uniform in-plane tension $N$ and two
damping mechanisms:
%
\begin{equation}
  \rho h\, \ddot w
  \;+\; c_1 \dot w
  \;+\; c_2\, \bilap \dot w
  \;+\; D\, \bilap w
  \;-\; N\, \lap w
  \;=\; f(\bm{x},t),
  \qquad
  D = \frac{E h^3}{12(1-\nu^2)},
  \label{eq:physical}
\end{equation}
%
where $\rho h$ is the surface density, $D$ the flexural rigidity, $\nu$ the
Poisson ratio ($\nu = 0.3$ in FemPlate), $c_1$ a viscous (air/support)
damping coefficient and $c_2$ a stiffness-proportional (material,
Kelvin--Voigt-type) damping coefficient. Dividing by $\rho h$ and rescaling
space and time so that the coefficients of $\ddot w$ and $\bilap w$ are
unity yields the \emph{scaled} equation actually used throughout the code:
%
\begin{equation}
  \ddot w + c_v\,\dot w + c_m\, \bilap\dot w + \bilap w - T\,\lap w = f ,
  \label{eq:scaled}
\end{equation}
%
with a single dimensionless \emph{tension-to-flexural-stiffness} parameter
$T \ge 0$. $T=0$ is a pure plate; $T \to \infty$ approaches the membrane
limit. Frequencies obtained from \eqref{eq:scaled} are dimensionless; the
plugin maps the first eigenfrequency onto the user parameter $f_1$ and all
others follow proportionally (\S\ref{sec:filterbank}), so the absolute
material constants never need to be specified.

\subsection{Boundary conditions}
\label{sec:bc}

The boundary $\partial\Omega$ is split by the user into segments, each
carrying one of four classical conditions:
%
\begin{center}
\begin{tabular}{llll}
\toprule
Name & Essential constraints & Natural conditions & DOFs constrained \\
\midrule
Free            & ---                          & $M_n=0$, $V_n=0$ & --- \\
Simple support  & $w=0$                        & $M_n=0$          & vertex $w$ \\
Clamp           & $w=0$, $\partial_n w=0$      & ---              & vertex $w$, edge $\partial_n w$ \\
Sliding support & $\partial_n w=0$             & $V_n=0$          & edge $\partial_n w$ \\
\bottomrule
\end{tabular}
\end{center}
%
where $M_n$ and $V_n$ are the bending moment and effective (Kirchhoff)
shear. Only the \emph{essential} conditions are enforced; the natural ones
are satisfied weakly by the variational formulation \cite{bathe}. The last
column anticipates \S\ref{sec:morley}: the four conditions map one-to-one
onto the degrees of freedom of the Morley element, which is the main reason
this element was chosen.

% ============================================================================
\section{Finite-element discretisation}
\label{sec:fem}

\subsection{Mesh generation}
\label{sec:mesh}

The library meshes an arbitrary closed polygon (the user's outline, sampled
from a closed Catmull--Rom spline \cite{catmull} for freehand shapes) with a
target element size $h$:
%
\begin{enumerate}
\item \emph{Boundary resampling.} The outline is walked at arc-length
  spacing $\approx h$. Vertices where the outline turns by more than
  $30^\circ$ are detected as corners and kept exactly, so rectangles remain
  rectangles. Every boundary sample stores its arc-length parameter
  $t\in[0,1)$, which is later used to attach the per-segment boundary
  conditions.
\item \emph{Interior lattice.} A hexagonal lattice of pitch $h$ fills the
  bounding box; points closer than $0.62\,h$ to the boundary or outside the
  polygon are discarded. A small deterministic jitter ($\pm 0.06\,h$)
  breaks the cocircular degeneracies that regular lattices would otherwise
  present to the triangulator.
\item \emph{Triangulation.} Bowyer--Watson incremental Delaunay insertion
  \cite{bowyer,watson}, with a super-triangle and cavity re-triangulation.
  The implementation is the plain $O(n^2)$ variant, entirely adequate for
  the $n \lesssim 2000$ points used here.
\item \emph{Clipping.} Triangles whose centroid falls outside the polygon
  are removed (this handles non-convex outlines), orientation is normalised
  to counter-clockwise, and an edge table with adjacency is built. Edges
  with a single adjacent triangle are boundary edges; they inherit the
  wrap-aware midpoint of their endpoints' arc parameters.
\end{enumerate}

\subsection{The Morley element}
\label{sec:morley}

The biharmonic operator requires $C^1$ continuity for conforming elements,
which leads to high-order elements (Argyris, Bell) that are expensive and
delicate. FemPlate instead uses the \emph{Morley triangle}
\cite{morley}, the simplest non-conforming plate bending element: on each
triangle the deflection is a full quadratic polynomial (6 coefficients)
determined by 6 degrees of freedom,
%
\begin{itemize}
\item the deflections $w_i$ at the three vertices,
\item the normal derivatives $(\partial_n w)_e$ at the three edge
  midpoints.
\end{itemize}
%
The element is non-conforming ($w$ is discontinuous along edges except at
vertices), yet it converges for fourth-order problems, and its eigenvalue
approximations converge at rate $O(h^2)$, \emph{from below}
\cite{rannacher,babuska-osborn} --- both properties are verified numerically
in \S\ref{sec:validation}. Crucially for FemPlate, its two DOF families are
exactly the quantities constrained by the four boundary conditions of
\S\ref{sec:bc}.

Rather than hard-coding the classical closed-form shape functions, each
element inverts a $6\times 6$ generalized Vandermonde system: writing the
local polynomial in centroid-centred monomials
$p(\bm{x}) = \sum_m c_m\, \mu_m(\bm{x})$,
$\{\mu_m\} = \{1, X, Y, X^2, XY, Y^2\}$, the matrix
$C_{im} = \ell_i(\mu_m)$ of the six DOF functionals $\ell_i$ is inverted by
Gaussian elimination with partial pivoting; column $j$ of $C^{-1}$ holds the
monomial coefficients of shape function $N_j$. Centroid-centred coordinates
keep $C$ well-conditioned at the element scale. The mid-edge normal is the
\emph{global} normal of the edge (stored with $v_0<v_1$, normal to the left
of $v_0\!\to\!v_1$), so the shared DOF has the same meaning in both adjacent
elements.

\subsection{Element matrices}

With constant curvatures per element (second derivatives of a quadratic),
the bending stiffness is exact:
%
\begin{equation}
  K^e_{ij} = A_e\, \bm\kappa_i^{\mathsf T} \mathbf{D}\, \bm\kappa_j,
  \qquad
  \bm\kappa_j = \begin{pmatrix} \partial_{xx} N_j \\ \partial_{yy} N_j \\
                 2\,\partial_{xy} N_j \end{pmatrix},
  \qquad
  \mathbf{D} = \begin{pmatrix} 1 & \nu & 0\\ \nu & 1 & 0\\
               0 & 0 & \tfrac{1-\nu}{2}\end{pmatrix},
\end{equation}
%
($A_e$ = element area; flexural rigidity is 1 after scaling). The tension
(geometric) matrix and the consistent mass matrix are integrated with
quadrature rules exact for their integrands:
%
\begin{align}
  G^e_{ij} &= \int_{e} \nabla N_i \cdot \nabla N_j \,\dd A
  &&\text{3-point midpoint rule (degree 2, exact),}\\
  M^e_{ij} &= \int_{e} N_i N_j \,\dd A
  &&\text{6-point degree-4 rule \cite{dunavant} (quartic, exact).}
\end{align}
%
$G$ is the discrete counterpart of $-\lap$ restricted to the plate space and
doubles as the \emph{geometric stiffness} familiar from buckling analysis
\cite{bathe}. Global assembly scatters $K^e, G^e, M^e$ into dense symmetric
matrices over the \emph{free} DOFs only: essential boundary conditions
(\S\ref{sec:bc}) are applied by elimination, using each boundary vertex's /
edge's arc parameter to look up its segment's condition. At a junction
between two segments the vertex takes the strongest condition
(clamp $>$ simple support $>$ sliding $>$ free).

\subsection{Validation}
\label{sec:validation}

The console test suite (\code{Tests/FemTests.cpp}) verifies the
implementation against classical analytical results \cite{leissa}:
%
\begin{itemize}
\item simply supported unit square: $\lambda_{mn} = \pi^4 (m^2+n^2)^2$,
  with the measured $O(h^2)$ convergence from below (error ratio $3.9$ when
  $h$ is halved) and the exact tension coefficient
  $g_{11} = 2\pi^2$ (\S\ref{sec:tension});
\item clamped circle: frequency ratios $2.081$, $3.414$, $3.893$ and the
  absolute fundamental $\omega_{01} = (\beta_{01}/R)^2$,
  $\beta_{01}^2 = 10.2158$ \cite{leissa}, within $2$--$5\%$ at moderate
  mesh density;
\item mixed clamped/free boundaries, mode-shape vanishing on the clamped
  side, and the first-order tension law against an exact re-solve.
\end{itemize}

% ============================================================================
\section{The eigenvalue problem}
\label{sec:eig}

\subsection{Formulation and solver}

Discretisation of \eqref{eq:scaled} (undamped part) yields the symmetric
generalized eigenproblem
%
\begin{equation}
  \left( K + T_0\, G \right) \bm\varphi_k \;=\; \lambda_k\, M \bm\varphi_k ,
  \qquad \lambda_k = \omega_k^2,
  \qquad \bm\varphi_k^{\mathsf T} M \bm\varphi_l = \delta_{kl},
  \label{eq:gep}
\end{equation}
%
solved at a \emph{reference tension} $T_0$ (the value of the Tension knob
when the user presses \emph{Compute}). With $n$ free DOFs
($n \approx$ vertices $+$ edges $\approx 4\times$ vertices, typically
$300$--$5000$ over the Grid range) and only the lowest $m \le 256$ modes
wanted (capped at $\sim n/6$ so the top modes stay resolved by the mesh),
the solver is
classical \emph{subspace iteration with Rayleigh--Ritz projection}
\cite{bathe-wilson,bathe}:
%
\begin{enumerate}
\item Form $A = K + T_0 G$ and the shifted operator $P = A + \sigma M$ with
  $\sigma = 10^{-5}\,\mathrm{tr}(A)/\mathrm{tr}(M)$. The shift keeps $P$
  positive definite when free (Neumann-type) boundaries leave rigid-body
  modes in $A$'s null space. $P$ is Cholesky-factored once ($O(n^3/3)$).
\item Iterate on a block of $p = m + \max(8, m/2)$ vectors:
  $Z \leftarrow P^{-1} M X$ (one inverse-power step towards the low
  spectrum), followed by Rayleigh--Ritz on $\mathrm{span}(Z)$: project
  $A_p = Z^{\mathsf T} A Z$, $M_p = Z^{\mathsf T} M Z$, solve the $p\times p$
  generalized problem by Cholesky reduction and cyclic Jacobi
  \cite{golub}, and take the Ritz vectors as the next block $X$.
\item Stop when the wanted eigenvalues are stationary: relative $10^{-6}$
  on the lower half, $10^{-4}$ on the upper half --- the slowest to
  converge, and the modes whose FEM discretisation error is orders of
  magnitude larger anyway (both thresholds are far below one cent in
  frequency); or after 60 iterations.
\end{enumerate}
%
Implementation notes that matter for speed: the iteration block is stored
\emph{transposed} (every iteration vector contiguous in memory, so all
$O(n^2 p)$ loops stream cache lines), the independent per-vector solves
and matrix-vector products are spread over a few worker threads, and the
Jacobi solver of the projected $p \times p$ problem uses a
\emph{relative} off-diagonal convergence test (the entries scale with
$\lambda^2 \sim 10^{12}$; an absolute threshold silently burns the full
sweep budget on every call). Together these give a $\sim\!4\times$
speed-up over the straightforward implementation ($23.6\,\mathrm{s}
\rightarrow 5.6\,\mathrm{s}$ for 128 modes on a 931-DOF mesh).
%
\paragraph{Sparse storage.}
A Morley element couples a degree of freedom only to the five others on
its own triangle, so an assembled row of $K$, $G$ or $M$ holds about
eleven non-zeros --- and, this being the useful part, about eleven
whatever the mesh density. Dense storage of the same matrix is $n$ per
row, so the fraction it wastes grows without bound: at $n = 5000$ the
three matrices are $24$ million doubles standing in for $55$ thousand real
numbers. They are therefore assembled in compressed sparse rows, in two
passes over the triangles: the first gathers nothing but the six DOF
indices of each element and builds the sparsity pattern (one pattern,
shared by all three matrices --- same mesh, same couplings), the second is
the expensive one and scatters element values into it. Both passes read
the DOF indices through the same helper, so the pattern cannot fail to
cover what the assembly scatters.

The shifted operator $P = A + \sigma M$ is still factorised \emph{densely},
which leaves one $n^2$ allocation and the $O(n^3/3)$ factorisation in
place. The working set is thus
%
\begin{equation}
  \underbrace{8 n^2}_{\text{Cholesky factor}}
  \;+\; \underbrace{32\,p\,n}_{\text{iteration block}}
  \;+\; \underbrace{\mathcal{O}(n)}_{\text{sparse operators}}
  \quad\text{bytes},
  \qquad\text{against } 32 n^2 \text{ dense},
  \label{eq:footprint}
\end{equation}
%
a factor of three to four in practice, with the matrix--vector products of
the iteration --- three per vector per sweep --- dropping from $O(n^2)$ to
$O(11n)$. Measured on the same machine, same meshes, both paths run
back to back:
%
\begin{center}
\begin{tabular}{rrrrrrr}
\toprule
& & \multicolumn{2}{c}{dense} & \multicolumn{2}{c}{sparse} & \\
\cmidrule(lr){3-4}\cmidrule(lr){5-6}
$n$ & modes & time & memory & time & memory & gain \\
\midrule
$2149$ & $64$  & $13.1$\,s  & $125$\,MB   & $7.3$\,s   & $36$\,MB  & $1.8\times$ / $3.4\times$ \\
$3841$ & $128$ & $114.8$\,s & $419$\,MB   & $59.9$\,s  & $122$\,MB & $1.9\times$ / $3.4\times$ \\
$6029$ & $128$ & $335.5$\,s & $1040$\,MB  & $171.0$\,s & $287$\,MB & $2.0\times$ / $3.6\times$ \\
\bottomrule
\end{tabular}
\end{center}
%
The two agree bit for bit, not merely to round-off:
compressed rows keep their column indices sorted, so a sparse row walk
sums the non-zeros in exactly the order a dense row walk does. The console
test suite asserts that equality, which makes it sharp enough to catch an
indexing slip.

\paragraph{What still sets the ceiling.}
The remaining $n^2$ is the dense factorisation, and it is what caps the
Grid knob. Over the Grid range the projected footprint is now $5$\,MB at
Grid 16 ($n=554$), $24$\,MB at 24 ($n=1243$), $63$\,MB at 32 ($n=2195$),
$132$\,MB at 40 ($n=3440$) and $243$\,MB at the top, Grid 48 ($n=4913$) ---
one third of what the same range cost dense. The editor shows the figure
next to the mesh size once it passes $150$\,MB. Removing the ceiling
altogether means factorising sparsely too: a bandwidth-reducing
(reverse Cuthill--McKee \cite{cuthill-mckee}) renumbering brings the
profile of these matrices down by more than an order of magnitude, and a
profile Cholesky over that envelope needs no symbolic analysis phase.
Measured on the actual DOF graph, that would take the Grid-48
factorisation from $39.5$ to $0.03$ Gflop and its storage from
$737$ to $5$\,MB. It is not, however, where the time goes: solving for 8
modes instead of 128 on the same mesh only drops the total from $85$ to
$18$\,s, so the factorisation is a fifth of the cost and the subspace
iteration is the rest. Sparsity would make the mesh nearly free and leave
the $O(n p^2)$ Rayleigh--Ritz projections --- which scale with the number
of modes asked for, not with the mesh --- as the thing to beat.

\paragraph{Code layout.}
None of the above is about plates, and it is not written as though it
were. The matrix storage, the Cholesky and Jacobi kernels and the subspace
iteration live in \code{FxmeTools/math}; \code{FxmeTools/acoustics} keeps
the boundary conditions, the Morley element and the modal quantities. The
eigensolver never sees a matrix at all --- it multiplies through a
\code{SymmetricOperator} and solves through an \code{SpdSolver} --- which
is what makes the storage a choice at the call site
(\code{ModalOptions::storage}, sparse by default, dense kept as the
reference the sparse path is validated against) rather than a rewrite, and
what will let a profile factorisation drop in as one more \code{SpdSolver}
with nothing else moving.

Finally, the computed modal data (eigenvalues, tension coefficients,
shapes, statistical tail) is cached inside the plugin state: loading a
session or preset republishes the model immediately, with no
recomputation. The mesh is not serialised --- \code{generateMesh} is a
pure deterministic function of the outline and density, so it is rebuilt
on load and checked against the stored vertex count.
%
Eigenvalues below an absolute threshold ($10^{-4}$; rigid-body and
near-rigid modes of free plates) are dropped. Eigenvectors are
mass-normalised, their vertex values are exported for display and for the
pointwise evaluations $\varphi_k(\bm{x})$ used by the synthesis (linear
interpolation on the located triangle; the mid-edge DOFs remain internal).
The whole computation runs on a background thread with progress reporting;
completed models are published to the audio thread by an atomic pointer
swap with generation-based reclamation.

\subsection{Audio-rate tension changes without re-solving}
\label{sec:tension}

The solver additionally returns, for every mode, the \emph{tension
sensitivity}
%
\begin{equation}
  g_k \;=\; \bm\varphi_k^{\mathsf T} G\, \bm\varphi_k
  \;\;\left(=\int_\Omega |\nabla \varphi_k|^2 \dd A \right) ,
\end{equation}
%
enabling the first-order frequency law used at audio rate:
%
\begin{equation}
  \boxed{\;\omega_k^2(T) \;=\; \lambda_k + (T - T_0)\, g_k\;}
  \label{eq:tensionlaw}
\end{equation}
%
This is legitimate because $T$ enters the operator \emph{linearly} and the
Rayleigh quotient is \emph{stationary} at eigenvectors: perturbing $T$
changes $\bm\varphi_k$ at first order in $(T-T_0)$, but that eigenvector
error affects the eigenvalue only at \emph{second} order. Equation
\eqref{eq:tensionlaw} is the standard eigenvalue derivative
$\partial\lambda_k/\partial T = \bm\varphi_k^{\mathsf T} G \bm\varphi_k$ of
structural sensitivity analysis \cite{fox-kapoor}, a special case of
classical perturbation theory \cite{courant,kato}; here it is exact at
$T_0$ with $O\!\left((T-T_0)^2\right)$ error. For the simply supported
square the law is exact for all $T$ (the eigenvectors do not depend on
$T$); in general, pressing \emph{Compute} re-anchors $T_0$ when the knob
has travelled far. The validation suite checks \eqref{eq:tensionlaw}
against an exact re-solve at $T=50$ (error $<1\%$).

Beyond the knob itself, \eqref{eq:tensionlaw} is the enabling mechanism for
the nonlinear glide of \S\ref{sec:berger}: \emph{dynamic} tension is just a
fast-moving $T$.

% ============================================================================
\section{Modal synthesis with resonant filter banks}
\label{sec:filterbank}

\subsection{Modal decomposition}

Expanding $w(\bm{x},t) = \sum_k q_k(t)\, \varphi_k(\bm{x})$ in
\eqref{eq:scaled} and using $M$-orthonormality decouples the dynamics into
damped oscillators
%
\begin{equation}
  \ddot q_k + 2\zeta_k \omega_k\, \dot q_k + \omega_k^2 q_k
  = \varphi_k(\bm{x}_h)\, F(t),
  \qquad
  w(\bm{x}_o, t) = \sum_k \varphi_k(\bm{x}_o)\, q_k(t),
  \label{eq:modal}
\end{equation}
%
for a point force $F$ at $\bm{x}_h$ and a pickup at $\bm{x}_o$ --- the
standard modal synthesis framework \cite{adrien,morrison,vandendoel,bilbao}.
FemPlate realises each oscillator as a second-order IIR resonator, exactly
as in the MechanOdd modal resonators: one constant-peak-gain band-pass
biquad per mode (RBJ ``cookbook'' coefficients \cite{rbj}) with
$Q_k = 1/(2\zeta_k)$, excited by $\varphi_k(\bm{x}_h) F$ and weighted on
output by $\varphi_k(\bm{x}_o)$. Up to 512 modes are active --- though the
bank stops at the highest one still below $0.47 f_s$, which at a high $f_1$
is a small fraction of them; the rest would cost a resonator each and
return silence. The bank is
allocation-free and retunable on the audio thread.

\subsection{Frequency and damping laws}

The user selects the fundamental in hertz and the bank plays at
%
\begin{equation}
  f_k \;=\; f_1\,
  \frac{\omega_k(T_{\mathrm{eff}})}{\omega_1(T_{\mathrm{knob}})} ,
  \qquad T_{\mathrm{eff}} = T_{\mathrm{knob}} + T_{\mathrm{dyn}} ,
  \label{eq:freqmap}
\end{equation}
%
with all $\omega_k$ given by \eqref{eq:tensionlaw}. The normalisation
reference matters: it is the fundamental at the \emph{static} knob tension,
not the instantaneous one. The static Tension knob therefore keeps mode 1
pinned at $f_1$ and only reshapes the overtone ratios (the timbre morphs
from plate towards membrane), whereas the \emph{dynamic} tension of
\S\ref{sec:berger} shifts the whole spectrum, fundamental included --- the
audible hardening glide. Normalising by the instantaneous fundamental
instead would cancel the glide of mode 1 and leave only the overtone-ratio
compression towards the membrane spectrum ($g_k/g_1$ ratios are lower than
the plate ratios $\sqrt{\lambda_k/\lambda_1}$), which is perceived as a
\emph{softening} --- the opposite of the physical behaviour. So the timbre
(inharmonicity pattern) is set by geometry, boundary conditions and
tension, while $f_1$ transposes it. Modes falling outside
$[20\,\mathrm{Hz},\ 0.47 f_s]$ are muted.

Projecting the two damping terms of \eqref{eq:scaled} onto mode $k$ gives
modal damping ratios (using $\bilap \varphi_k \approx \omega_k^2 \varphi_k$
for the material term)
%
\begin{equation}
  \zeta_k \;=\; \underbrace{\frac{c_v}{2\omega_k}}_{\text{viscous}}
          \;+\; \underbrace{\frac{c_m\, \omega_k}{2}}_{\text{material}}
  \;\;=\;\; \frac{\zeta_v}{\nu_k} + \zeta_m\, \nu_k ,
  \label{eq:damping}
\end{equation}
%
where the knobs $\zeta_v, \zeta_m$ are simply the two contributions to the
damping ratio \emph{of mode 1}. Viscous damping therefore lets high modes
ring relatively longer (constant absolute bandwidth), material damping
kills them faster --- the two classical, perceptually distinct decay
profiles of struck plates \cite{chaigne,fletcher-rossing}.

\subsection{Excitation, pickup, gain compensation}

A mouse hit at $\bm{x}_h$ launches a half-sine force pulse (a standard
hammer/impact model \cite{chaigne})
%
\begin{equation}
  F(t) = F_0 \sin\!\left(\pi t/\tau\right), \qquad 0 \le t \le \tau,
\end{equation}
%
with duration $\tau$ (\emph{Hammer} knob, 0.1--50\,ms: short = bright,
long = mellow, since the pulse spectrum rolls off above $\approx 1/\tau$)
and amplitude $F_0$ (\emph{Force} knob $\times$ click velocity). Up to 8
hammers overlap polyphonically, each carrying its own weight vector
$\varphi_k(\bm{x}_h)$. The external audio input is injected identically at
the last hit point (\emph{effect mode}), and the output point
$\bm{x}_o$ is a pair of automatable parameters.

Because a constant-peak band-pass transmits energy proportionally to its
bandwidth ($\propto \zeta_k$), lightly damped modes would ring much quieter
than damped ones. Each mode is therefore compensated by
$c_k = \sqrt{\zeta_{\mathrm{ref}}/\zeta_k}$ (clamped to $[0.5, 32]$,
$\zeta_{\mathrm{ref}}=10^{-2}$), the loudness compensation inherited from
the MechanOdd modal resonators:
%
\begin{equation}
  c_k \;=\; \sqrt{\zeta_{\mathrm{ref}}/\zeta_k} .
  \label{eq:compensation}
\end{equation}
%
It equalises each mode's ring energy rather than its peak, which has a
consequence for anything driven from the output signal --- see
\S\ref{par:drive}. Retunes triggered by parameter changes
rewrite only the biquad coefficients (the states are preserved), which is
inaudible for the small throttled steps used here.

% ============================================================================
\section{Nonlinear approximation}
\label{sec:nonlinear}

Large-amplitude plate vibration is governed by the von K\'arm\'an
equations \cite{vonkarman,chu-herrmann}: transverse deflection stretches
the mid-plane, producing an amplitude-dependent membrane stress that (i)
raises all frequencies (the gong/tom \emph{pitch glide}) and (ii) couples
the modes, cascading energy towards high frequencies (the cymbal
\emph{crash}, up to wave turbulence \cite{touze-turbulence}). Exact modal
von K\'arm\'an synthesis requires the cubic coupling tensors
$\Gamma_{klmn}$ --- $O(N^3)$--$O(N^4)$ storage and work
\cite{ducceschi-touze,bilbao-vk} --- which is beyond a realtime budget at
a few hundred modes. FemPlate implements two lightweight surrogates, one for each
phenomenon; \S\ref{sec:related} compares them with the direct simulation
of the full system, which is affordable in real time on a rectangle.

\subsection{Dynamic tension: the Berger approximation}
\label{sec:berger}

Berger \cite{berger} observed that for many plate problems the von
K\'arm\'an membrane coupling is well approximated by a \emph{spatially
uniform} tension proportional to the integrated stretching:
%
\begin{equation}
  \ddot w + \bilap w
  - \Bigl( T_{\mathrm{knob}}
     + \underbrace{\kappa \int_\Omega |\nabla w|^2 \dd A}_{T_{\mathrm{dyn}}}
    \Bigr) \lap w = f .
  \label{eq:berger}
\end{equation}
%
In modal coordinates
$\int |\nabla w|^2 = \sum_{kl} q_k q_l\, \bm\varphi_k^{\mathsf T} G
\bm\varphi_l \approx \sum_k g_k q_k^2$: the stretching integral is measured
by the very same $G$ and $g_k$ of \S\ref{sec:tension}, and the resulting
uniform $T_{\mathrm{dyn}}$ feeds straight into the retune law
\eqref{eq:tensionlaw}. The Berger model captures the hardening
(glide-up-and-relax) behaviour of plates and membranes and is a standard
efficient approximation in the nonlinear vibration literature
\cite{berger,nayfeh}; its perceptual signature on percussion is documented
by Legge \& Fletcher and Rossing \cite{legge-fletcher,fletcher-rossing}.

\paragraph{Amplitude estimation and calibration.}
\label{par:driver}
The stretching integral needs no field reconstruction: the modes are
mass-orthonormal, so $\int_\Omega w^2\,\dd A = \sum_k q_k^2$ exactly and
$\int_\Omega |\nabla w|^2 \dd A \approx \sum_k g_k q_k^2$, one multiply-add
per mode and per sample. What the filter bank holds is not $q_k$, however.
A $0$\,dB-peak band-pass fed the true modal force is exactly
$y_k = 2\zeta_k\omega_k\,\dot q_k$ --- a \emph{velocity}, normalised --- so
%
\begin{equation}
  q_k \;=\; \frac{y_k}{2\zeta_k\omega_k^2},
  \qquad
  \gamma \;=\; \min\!\Bigl(\gamma_{\max},\;
        \alpha\, \cdot \mathrm{nonlin} \cdot
        \Bigl\langle \sum_k W_k\, y_k^2 \Bigr\rangle \Bigr),
  \qquad
  W_k = \frac{g_k}{g_1}
        \left(\frac{\zeta_{\mathrm{ref}}}{\zeta_k}\right)^{\!2}
        \frac{1}{\nu_k^4},
  \label{eq:gamma}
\end{equation}
%
with $T_{\mathrm{dyn}} = \gamma\, \omega_1^2(T_{\mathrm{knob}})/g_1$,
$\alpha = 500$, $\gamma_{\max}=4$, and $\langle\cdot\rangle$ a 20-ms
one-pole envelope. The missing $1/(4\zeta^2\omega^4)$ --- ten decades
across the bank --- is why a naive $\sum_k g_k\langle y_k^2\rangle$
underestimates the stretching by orders of magnitude. Being derived from
the velocity, each mode is a quarter period ahead of its true
displacement; the cross terms $k \ne l$ and that phase both average out in
the envelope.

By construction $\omega_1^2 \mapsto \omega_1^2\,(1+\gamma)$ exactly, and
every other mode follows its own physical sensitivity $g_k$: $\gamma$ is
plate-independent ($\gamma = 0.4$ always means ``mode 1 up $\approx 3$
semitones''). It is also \emph{damping}-independent by construction: after
an impulse $y_k \sim 2\zeta_k\omega_k$, so $\zeta_k$ cancels in
$W_k y_k^2$ and the peak driver varies by less than $12\%$ over
$\zeta_v, \zeta_m \in [10^{-5}, 10^{-3}]$ on the default plate. And being
a quantity integrated over $\Omega$, it is the same wherever one listens:
the glide belongs to the plate, not to the pickup.

\paragraph{Realtime scheduling.}
$\gamma$ is updated every 32 samples, slewed by a factor $0.2$ per update
($\approx$ ms-scale smoothing), and the biquads are rewritten only when
mode 1 would move by more than $\approx 2$ cents
($|\Delta\gamma| > 0.0024\,(1+\gamma)$, since
$\Delta\omega/\omega = \tfrac12 \Delta\gamma/(1+\gamma)$). Coefficient
rewrites preserve filter states; the throttling keeps them free of zipper
artefacts, following the same strategy as MechanOdd's gate-morphed damping.

\subsection{Mode cascade: a multi-band cubic ladder}
\label{sec:cascade}

A memoryless cubic of a signal band-limited to $[0, B_\omega]$ produces
inter-modulation products only up to $3 B_\omega$, concentrated near the
original band. A single ``distort the low modes, inject into the high
ones'' pass therefore \emph{cannot} reach the top of the bank: with
sources below $\sim\!12 f_1$ the products die out by $\sim\!3$\,kHz while
the bank extends to $9$\,kHz, and all that remains audible of such a
scheme is the side effect of its overlap floor --- high-frequency damping.
The physical cascade climbs the spectrum \emph{iteratively}: products
beget products, stage after stage \cite{touze-turbulence}.

FemPlate reproduces this with a \emph{windowed} band ladder. The active
modes are split into $N_b = 8$ contiguous index bands
$\mathcal{B}_0 \ldots \mathcal{B}_7$ (equal index ranges are equal
frequency ranges, plate modal density being constant); band $b$ is driven
by the cubic of the $W$ bands directly below it only (default $W = 4$),
%
\begin{equation}
  w^{(b)}(t) = \frac{S}{\sqrt{|\Omega|}} \sum_{c = b-W}^{b-1}
              \frac{1}{n_c} \sum_{k \in \mathcal{B}_c} v_k\, y_k(t),
  \qquad
  F_b(t) \;=\; A\, g_b(t)\, \tanh^3\!\bigl(B\, w^{(b)}(t)\bigr),
  \label{eq:cascade}
\end{equation}
%
fed to mode $l \in \mathcal{B}_b$ with weight $\varphi_l(\bm{x}_h)/c_l$
(dividing by the bandwidth compensation $c_l$, which is re-applied on
output, makes the audible cascade level independent of the damping
settings). The source is the plate's \emph{motion}, over the whole plate:
$v_k = \zeta_{\mathrm{ref}}/(\zeta_k \nu_k)$ is the modal velocity scale,
$n_c = \sum_{k \in \mathcal{B}_c} 1/\nu_k$ normalises each band,
$1/\sqrt{|\Omega|}$ is what a mass-normalised $\varphi_k$ is typically
worth, and $S = 3$ is the one global gain (\S\ref{par:drive}). The source \emph{window} matters: summing all lower bands
would let the loud strike band drive the top bands quasi-directly and the
whole spectrum lights up within milliseconds; with the window, each stage
extends the spectral reach by at most a factor 3 and energy genuinely
climbs rung by rung. The per-band gate $g_b(t) \in [0,1]$ adds
\emph{transfer inertia}: it follows the source presence through an
attack/release envelope whose attack time grows with the band's height on
the ladder ($\tau_b = 30\,\mathrm{ms} \times b$ and a $2\,\mathrm{s}$
release in the voiced defaults), so the brightening is a progressive glow
and the pumping tails off smoothly instead of collapsing with the cube of
the decaying ring. Voicing found a hard-saturating transfer most
convincing: a large drive ($B = 16$) with a small injection gain
($A = 1.1$, capped at $1.5$ --- beyond that the level becomes
unpleasant). A
cubic is the leading nonlinearity of von K\'arm\'an dynamics
\cite{ducceschi-touze}; at low amplitude \eqref{eq:cascade} vanishes as
$w^3$ (the linear model is recovered exactly), and at high amplitude the
bands brighten in sequence, reproducing the amplitude-gated crash of
cymbals and gongs \cite{fletcher-rossing}. All cascade internals ($A$,
$B$, the source window, the gate times and the overlap floor below) are
exposed as plugin parameters while the effect is being voiced; the values
quoted here are the defaults. Note also that the material damping law
$\zeta \propto \nu$ dominates the high-mode decay when hunting long
shimmer: at $\zeta_m = 10^{-4}$ a mode at $\nu = 80$ decays in
milliseconds whatever the cascade does.

\paragraph{Modal-overlap floor.}
The cubic's inter-modulation products are broadband, but at typical plate
dampings the receiving resonances are needles: bandwidth $2\zeta_l f_l$ of
a few hertz against a mean mode spacing $\Delta f \approx 0.64 f_1$ (the
plate Weyl law gives a constant modal density $\dd N/\dd\nu = \pi/2$ for
the unit square), i.e.\ a modal overlap of order $0.1$. Most of the
injected energy then falls \emph{between} modes and is rejected --- heard
as distortion of isolated partials rather than a cascade, whereas the
crash regime of real cymbals has overlap $\gg 1$ \cite{fletcher-rossing}.
The target bandwidths are therefore floored, proportionally to the cascade
control, at an adjustable fraction (\emph{Overlap}, default $0.1$) of the
\emph{local} mode spacing (measured on the actual $f_k$ array),
%
\begin{equation}
  \zeta_l \;\ge\; \mathrm{cascade} \times \mathrm{Overlap} \times
      \frac{\tfrac12 (f_{l+1} - f_{l-1})}{2 f_l},
  \qquad l \notin \mathcal{B}_0,
\end{equation}
%
turning the receiving comb into a quasi-continuum when the cascade is
engaged (and, physically consistently, shortening the high modes' ring ---
nonlinear energy spreading is itself a loss channel for individual modes).
At $\mathrm{cascade}=0$ the linear model is untouched.

\paragraph{Why the source is the motion and not the sound.}
\label{par:drive}
The obvious source signal for \eqref{eq:cascade} is the one already being
summed for the output, $\varphi_k(\bm{x}_o)\,c_k\,y_k$ --- and that is
precisely what it must not be. Two things stand between the audible signal
and the plate's motion, and the ladder's cubic would raise both to the
third power.

The first is the loudness compensation \eqref{eq:compensation}. A band-pass
ring has peak $2\zeta_k\omega_k$, so the audible peak of a struck mode is
%
\begin{equation}
  c_k \cdot 2\zeta_k \omega_k \;=\; 2\omega_k
      \sqrt{\zeta_{\mathrm{ref}}\,\zeta_k} \;\propto\; \sqrt{\zeta_k} :
  \label{eq:peaklaw}
\end{equation}
%
$c_k$ equalises each mode's ring \emph{energy} by trading duration for
level, so in the audible signal a more damped mode peaks \emph{higher}.
That is sound design, and it is the right choice for what reaches the ear,
but it is not what the plate does: a struck plate's modal amplitude is set
by the impulse, and damping can only take energy away. Weighting the
source by $v_k \propto 1/(\zeta_k\nu_k)$ removes the damping entirely,
since $y_k \propto \zeta_k$; what is left is the modal velocity.

The per-band divisor $n_c$ then supplies the \emph{frequency} balance the
ladder needs, which \eqref{eq:peaklaw} had been providing as a side
effect: physical modal amplitudes fall steeply with frequency, and a raw
velocity sum outweighs band 1 with band 0 by a factor $\sim\!900$,
starving the upper rungs. $n_c$ is fixed by the mode indices alone, so no
damping re-enters.

The second is the pickup. A mode nodal at $\bm{x}_o$ contributes nothing
to the audible signal however hard the plate is moving there, and a real
plate's cubic coupling does not know where the microphone is. The source
weight is therefore the constant $1/\sqrt{|\Omega|}$, the typical size of
a mass-normalised mode shape ($\int_\Omega \varphi_k^2 \dd A = 1$), which
holds the drive's level on any geometry while tying it to no listening
position: across five pickup positions the high-mode motion the cascade
creates is identical to the last digit. What the pickup changes is what
one \emph{hears} of that motion, which is what a pickup is for. The
\emph{injection} does keep $\varphi_l(\bm{x}_h)$: the hit point is where
the plate deflects most, the one position the nonlinearity has a physical
claim to, and it is what ties the cascade's character to how the plate is
played.

The Cascade knob therefore acts directly, with no damping correction on
it. \code{Tests/CascadeMeasure.cpp} renders one strike through the real
synthesis and asserts what the construction promises --- the drive is
damping-free, so any real spread would mean a damping term had crept back
into it:
%
\begin{center}
\begin{tabular}{lc}
\toprule
Shimmer at 300\,ms, spread over & dB \\
\midrule
Viscous $\zeta_v \in [10^{-5}, 10^{-3}]$              & $5.0$ \\
Material $\zeta_m \in [7\cdot10^{-6}, 3\cdot10^{-4}]$ & $6.6$ \\
Output point, five positions (drive, measured on the modes) & $0.00$ \\
\bottomrule
\end{tabular}
\end{center}
%
The few dB left on the damping rows are not drive: the targets ring for
longer or shorter with the damping, and the middle rungs are sources for
the upper ones. Meanwhile the Cascade knob spans $18$\,dB over its travel
and a force-$10$ hit shimmers $45$\,dB more than a force-$1$ one --- the
effect stays amplitude-gated, which is the part that \emph{is} physical.

\paragraph{Depletion (low-frequency loss).}
The transfer \eqref{eq:cascade} is additive: it copies energy upward but
takes none from the sources, so the struck low modes ring on untouched
under the shimmer --- audibly two disconnected layers. In the real
cascade the low modes \emph{lose} the energy they hand upward. FemPlate
models this as an extra dissipation on the source bands: each sample, the
filter states of the modes in band $b$ are scaled by
%
\begin{equation}
  d_b \;=\; 1 - \mathrm{Deplete} \times \mathrm{cascade} \times
            \varepsilon_0\, \bar g_{b}^{\,\uparrow},
  \qquad \varepsilon_0 = 3\cdot 10^{-4}
  \quad (\mathrm{Deplete} = 0.07 \text{ by default}),
\end{equation}
%
where $\bar g_b^{\,\uparrow}$ is the mean gate level of the bands fed by
band $b$: the lows decay exactly when, and as strongly as, they are
pumping the highs, gluing the two ends of the spectrum. The mechanism is
purely dissipative (it only removes energy), so it cannot affect
stability.

\paragraph{Stability by structure.}
The obvious way to spread energy upward is a feedback loop closed through
the bank's own output, and it does not survive contact with the gains
involved. Around mode $k$'s resonance such a loop transmits with
$|\varphi_k(\bm{x}_h)|\cdot|\varphi_k(\bm{x}_o)|\cdot c_k$, and the
compensation $c_k \le 32$ must not be forgotten: an unnormalised loop is
super-critical and latches into a full-scale limit cycle the moment the
cubic's slope wakes up. Bounding the loop slope by the small-gain theorem
\cite{zames,khalil} restores stability but is so conservative (the bound
assumes worst-case coherent feedback at the single worst mode) that the
effect becomes inaudible. The ladder \eqref{eq:cascade} dissolves the
dilemma structurally: band $b$ receives only from bands $c<b$ and
contributes only to bands $c>b$, so the transfer graph over bands is a
strict DAG --- no loop exists at any stage, stability is unconditional,
and the injection gain $A$ is free to be set for audibility
($A = \mathrm{Amp} \times \mathrm{cascade}$, the $\tanh^3$ keeping every
stage bounded).

\subsection{The statistical mode tail}
\label{sec:tail}

Even a 256-mode FEM bank remains a sparse comb at the top of the
spectrum, and solving for many hundreds of accurate FEM modes is
prohibitive (\S\ref{sec:eig}) --- while a convincing crash needs a dense
receiving continuum up to Nyquist. FemPlate therefore fills the bank
(to 512 modes in all) with \emph{synthetic} modes above the FEM set,
built from asymptotics instead of the discrete operator:
%
\begin{itemize}
\item \emph{Frequencies}: the plate Weyl law gives a constant modal
  density in $\omega$; the tail continues at the spacing
  $\Delta\omega$ measured on the top half of the FEM spectrum, with
  gaps jittered by $\pm 40\%$ (deterministically seeded) to avoid
  periodicity artefacts.
\item \emph{Tension sensitivities}: for an asymptotic mode of wavenumber
  $\kappa$, $\lambda = \kappa^4$ and $g = \kappa^2$, hence
  $g = \sqrt{\lambda}$ exactly (verified by the simply supported square,
  where $g_{mn} = \pi^2(m^2{+}n^2) = \sqrt{\lambda_{mn}}$) --- so the
  tail obeys the tension law \eqref{eq:tensionlaw} and the Berger glide
  like any FEM mode.
\item \emph{Shapes}: high-frequency eigenfunctions of irregular domains
  behave like superpositions of random plane waves (Berry's random-wave
  model \cite{berry}); each tail mode carries one,
  $\varphi(\bm{x}) = \sqrt{2/A}\, \cos(\bm{\kappa}\cdot\bm{x} + \psi)$
  with $|\bm{\kappa}| = \lambda^{1/4}$, random direction and phase, and
  the $\sqrt{2/A}$ prefactor reproducing the mass normalisation
  ($\overline{\varphi^2} = 1/A$). Hit and pickup positions therefore
  still modulate the tail response.
\end{itemize}
%
The tail modes enter the synthesis exactly like FEM modes (hammer
excitation, damping law, cascade bands, dynamic tension); only the
contour display is FEM-only, since tail modes have no mesh shape. The
identity-defining low modes stay exact, and the top of the bank becomes
the quasi-continuum the cascade ladder needs --- the reasoning is that of
statistical energy analysis \cite{lyon}: above a few dozen mode indices,
individual eigenfrequencies are perceptually interchangeable and only
their density and statistics matter.

% ============================================================================
\subsection{Relation to other nonlinear plate synthesis}
\label{sec:related}

The closest published work to FemPlate's objective is the real-time gong
synthesis of Bilbao, Webb, Wang and Ducceschi \cite{bilbao-gong23}, which
reaches the same sounds --- pitch glide, crash, swell --- from the opposite
direction, and the contrast is worth stating plainly because it delimits
what the surrogates of \S\ref{sec:berger} and \S\ref{sec:cascade} can and
cannot claim.

Direct time-domain simulation of the von K\'arm\'an plate for synthesis
goes back to Bilbao's offline work \cite{bilbao-nlplates,bilbao-vk}; what
\cite{bilbao-gong23} adds is the algorithmic work needed to bring it inside
a real-time budget.

That work simulates the dynamic F\"oppl--von K\'arm\'an system directly: the
deflection $w$ and the Airy stress function $\Phi$ are stepped on a uniform
Cartesian grid, with the nonlinearity made fully explicit by a scalar
auxiliary variable (energy quadratisation), a rank-one system solved in
closed form by Sherman--Morrison, and a discrete energy balance that yields
a provable stability condition \cite{bilbao-sav,bilbao-gong23}. The
nonlinear effects are not modelled separately: they are consequences of the
one equation, so the glides come out non-uniform across partials and the
crash appears on its own as the amplitude rises. The price is paid in three
places:

\begin{itemize}
\item \emph{Geometry.} The domain is a rectangle with simply supported
  conditions on both $w$ and $\Phi$. This is not incidental. The run-time
  bottleneck is the biharmonic solve $\Delta\Delta\Phi = -L(w,w)$ needed at
  every sample, which consumes $72$--$76\%$ of the time step, and the fast
  solver that makes it affordable \cite{wang-puckette} exploits the
  Toeplitz-block-Toeplitz structure that the rectangle and those boundary
  conditions produce. An arbitrary outline would cost the solver, and with
  it the real-time budget.
\item \emph{Cost.} Reported timings are $0.15$--$0.57$\,s of computation per
  second of audio on a recent laptop CPU, and $0.26$--$0.99$\,s on an older
  one, for grids of roughly $300$ to $720$ interior points. FemPlate's full
  bank --- 256
  modes with both the Berger glide and the cascade engaged --- measures
  $0.054$\,s per second of audio on an Intel i7-10810U, roughly an order of
  magnitude less, because nothing is solved at run time: the eigenproblem
  was solved once, offline (\S\ref{sec:eig}; $0.5$\,s for the default
  plate), and what remains is a bank of biquads whose coefficients are
  rewritten when the tension moves.
\item \emph{Resolution coupling.} The grid spacing is tied to the sample
  rate by the stability condition, so spatial resolution is not a free
  modelling choice. In a modal scheme the mesh density and the audio rate
  are independent.
\end{itemize}

What is given up in exchange is exactness, and that paper names the
limitation precisely: Berger's approximation ``[is] able to replicate pitch
glides, but not the energy cascade to high frequencies characteristic of
crashes'' \cite{bilbao-gong23}. That is correct, and it is why
\S\ref{sec:cascade} exists as a \emph{separate} mechanism rather than
falling out of \S\ref{sec:berger}. Two consequences follow. First, the
glide here is driven by a single scalar $T_{\mathrm{dyn}}$: every mode
follows it through its own $g_k$, so the pattern of glides is fixed by
geometry and boundary conditions, where in a von K\'arm\'an solution mode
$k$'s frequency shift depends on which modes actually carry the energy and
therefore evolves through the decay. Second, the cascade of
\S\ref{sec:cascade} is a caricature of the coupling, not the coupling: it
reproduces the phenomenon --- energy climbing the spectrum rung by rung,
gated by amplitude --- with a band ladder whose gains are voiced by ear,
and it carries no energy balance. Stability comes from the transfer graph
being acyclic rather than from a numerical invariant.

The exact modal route sits between the two. Projecting von K\'arm\'an onto
the linear modes gives cubic coupling tensors $\Gamma_{klmn}$, which is the
basis of the gong and cymbal synthesis of Ducceschi and Touz\'e
\cite{ducceschi-touze} and of the analytical work on circular plates and
shallow shells by Touz\'e, Thomas and Chaigne
\cite{touze-circular,thomas-shells}; Nguyen and Touz\'e extend it to plates
of variable thickness for cymbals \cite{nguyen-touze}. That route keeps the
modal picture --- and would keep FemPlate's arbitrary geometry, since the
tensors can in principle be assembled from any set of FEM modes --- but the
$O(N^3)$--$O(N^4)$ tensor is what rules it out at $N = 256$ in a plugin,
and it is the reason the two cheap surrogates were chosen instead.

Finally, the mechanism of \S\ref{sec:berger} is not new in synthesis: an
amplitude-dependent tension computed from a global stretching estimate, fed
back into the linear model's tuning, is the plate analogue of the
tension-modulation models used for plucked strings, where the same
first-order retuning reproduces the pitch drop of a hard pluck
\cite{tolonen}. What is specific here is that the tension enters through
the FEM sensitivities $g_k$ of \S\ref{sec:tension}, so it costs one scalar
per audio block on any drawn shape.

\section{Implementation summary}

\begin{center}
\begin{tabular}{llp{7.2cm}}
\toprule
Stage & Thread & Notes \\
\midrule
Shape / boundary editing & message & \code{ShapeCanvas}; outline + arc-parameterised segments \\
Mesh generation & message & milliseconds; preview in \code{FemViewComponent} \\
Eigensolution \eqref{eq:gep} & background & \code{BackgroundTaskRunner}; progress; seconds \\
Model publication & message $\to$ audio & atomic pointer, generation-acknowledged reclamation \\
Filter bank \eqref{eq:modal} & audio & allocation-free; retunes throttled \\
Nonlinearities \eqref{eq:gamma}, \eqref{eq:cascade} & audio & decimated $\gamma$ update; bounded feedback \\
\bottomrule
\end{tabular}
\end{center}

The numerical core is pure C++17 with no JUCE dependency and is validated
by a console test suite; it lives in the \code{FxmeTools} submodule, split
between \code{FxmeTools/acoustics} (\code{FemMesh}, \code{PlateModes}: the
plate physics) and \code{FxmeTools/math} (matrix storage, Cholesky,
subspace eigensolver: the part that has never heard of a plate), so that
either half is reusable on its own.

% ============================================================================
\begin{thebibliography}{99}

\bibitem{leissa} A.~W. Leissa, \emph{Vibration of Plates}, NASA SP-160,
  Washington DC, 1969.

\bibitem{timoshenko} S.~Timoshenko, S.~Woinowsky-Krieger, \emph{Theory of
  Plates and Shells}, 2nd ed., McGraw-Hill, 1959.

\bibitem{bathe} K.-J. Bathe, \emph{Finite Element Procedures}, 2nd ed.,
  Prentice Hall / K.-J. Bathe, 2014.

\bibitem{catmull} E.~Catmull, R.~Rom, ``A class of local interpolating
  splines'', in \emph{Computer Aided Geometric Design}, Academic Press,
  1974, pp.~317--326.

\bibitem{bowyer} A.~Bowyer, ``Computing Dirichlet tessellations'',
  \emph{The Computer Journal} 24(2), 1981, 162--166.

\bibitem{watson} D.~F. Watson, ``Computing the $n$-dimensional Delaunay
  tessellation with application to Vorono\"i polytopes'', \emph{The
  Computer Journal} 24(2), 1981, 167--172.

\bibitem{morley} L.~S.~D. Morley, ``The triangular equilibrium element in
  the solution of plate bending problems'', \emph{Aeronautical Quarterly}
  19, 1968, 149--169.

\bibitem{rannacher} R.~Rannacher, ``Nonconforming finite element methods
  for eigenvalue problems in linear plate theory'', \emph{Numerische
  Mathematik} 33, 1979, 23--42.

\bibitem{babuska-osborn} I.~Babu\v{s}ka, J.~Osborn, ``Eigenvalue
  problems'', in \emph{Handbook of Numerical Analysis}, Vol.~II,
  North-Holland, 1991, 641--787.

\bibitem{cuthill-mckee} E.~Cuthill, J.~McKee, ``Reducing the bandwidth of
  sparse symmetric matrices'', in \emph{Proc. 24th National Conference of
  the ACM}, 1969, pp.~157--172.

\bibitem{dunavant} D.~A. Dunavant, ``High degree efficient symmetrical
  Gaussian quadrature rules for the triangle'', \emph{International
  Journal for Numerical Methods in Engineering} 21, 1985, 1129--1148.

\bibitem{bathe-wilson} K.-J. Bathe, E.~L. Wilson, ``Solution methods for
  eigenvalue problems in structural mechanics'', \emph{International
  Journal for Numerical Methods in Engineering} 6, 1973, 213--226.

\bibitem{golub} G.~H. Golub, C.~F. Van Loan, \emph{Matrix Computations},
  4th ed., Johns Hopkins University Press, 2013.

\bibitem{courant} R.~Courant, D.~Hilbert, \emph{Methods of Mathematical
  Physics}, Vol.~I, Interscience, 1953.

\bibitem{kato} T.~Kato, \emph{Perturbation Theory for Linear Operators},
  2nd ed., Springer, 1976.

\bibitem{fox-kapoor} R.~L. Fox, M.~P. Kapoor, ``Rates of change of
  eigenvalues and eigenvectors'', \emph{AIAA Journal} 6(12), 1968,
  2426--2429.

\bibitem{adrien} J.-M. Adrien, ``The missing link: modal synthesis'', in
  G.~De Poli, A.~Piccialli, C.~Roads (eds.), \emph{Representations of
  Musical Signals}, MIT Press, 1991, 269--298.

\bibitem{morrison} J.~D. Morrison, J.-M. Adrien, ``MOSAIC: a framework for
  modal synthesis'', \emph{Computer Music Journal} 17(1), 1993, 45--56.

\bibitem{vandendoel} K.~van den Doel, D.~K. Pai, ``Modal synthesis for
  vibrating objects'', in K.~Greenebaum (ed.), \emph{Audio Anecdotes III},
  A.~K. Peters, 2003.

\bibitem{bilbao} S.~Bilbao, \emph{Numerical Sound Synthesis: Finite
  Difference Schemes and Simulation in Musical Acoustics}, Wiley, 2009.

\bibitem{rbj} R.~Bristow-Johnson, ``Cookbook formulae for audio EQ biquad
  filter coefficients'' (the ``Audio EQ Cookbook''),
  \url{https://www.w3.org/TR/audio-eq-cookbook/}.

\bibitem{chaigne} A.~Chaigne, C.~Lambourg, ``Time-domain simulation of
  damped impacted plates. I. Theory and experiments'', \emph{Journal of
  the Acoustical Society of America} 109(4), 2001, 1422--1432.

\bibitem{fletcher-rossing} N.~H. Fletcher, T.~D. Rossing, \emph{The
  Physics of Musical Instruments}, 2nd ed., Springer, 1998.

\bibitem{vonkarman} T.~von K\'arm\'an, ``Festigkeitsprobleme im
  Maschinenbau'', \emph{Encyklop\"adie der mathematischen Wissenschaften}
  IV/4, 1910, 311--385.

\bibitem{chu-herrmann} H.-N. Chu, G.~Herrmann, ``Influence of large
  amplitudes on free flexural vibrations of rectangular elastic plates'',
  \emph{Journal of Applied Mechanics} 23, 1956, 532--540.

\bibitem{berger} H.~M. Berger, ``A new approach to the analysis of large
  deflections of plates'', \emph{Journal of Applied Mechanics} 22, 1955,
  465--472.

\bibitem{nayfeh} A.~H. Nayfeh, D.~T. Mook, \emph{Nonlinear Oscillations},
  Wiley, 1979.

\bibitem{legge-fletcher} K.~A. Legge, N.~H. Fletcher, ``Nonlinearity,
  chaos, and the sound of shallow gongs'', \emph{Journal of the Acoustical
  Society of America} 86(6), 1989, 2439--2443.

\bibitem{ducceschi-touze} M.~Ducceschi, C.~Touz\'e, ``Modal approach for
  nonlinear vibrations of damped impacted plates: application to sound
  synthesis of gongs and cymbals'', \emph{Journal of Sound and Vibration}
  344, 2015, 313--331.

\bibitem{bilbao-gong23} S.~Bilbao, C.~Webb, Z.~Wang, M.~Ducceschi,
  ``Real-time gong synthesis'', in \emph{Proceedings of the 26th
  International Conference on Digital Audio Effects (DAFx23)}, Copenhagen,
  2023, pp.~1--8.

\bibitem{bilbao-nlplates} S.~Bilbao, ``Sound synthesis for nonlinear
  plates'', in \emph{Proceedings of the 8th International Conference on
  Digital Audio Effects (DAFx-05)}, Madrid, 2005, pp.~243--248.

\bibitem{bilbao-sav} S.~Bilbao, M.~Ducceschi, F.~Zama, ``Explicit exactly
  energy-conserving methods for Hamiltonian systems'', \emph{Journal of
  Computational Physics}, 2023.

\bibitem{wang-puckette} Z.~Wang, M.~Puckette, ``A fast algorithm for the
  inversion of the biharmonic in plate dynamics applications'', in
  \emph{Proceedings of the 5th Stockholm Music Acoustics Conference},
  Stockholm, 2023.

\bibitem{touze-circular} C.~Touz\'e, O.~Thomas, A.~Chaigne, ``Asymmetric
  non-linear forced vibrations of free-edge circular plates, part I:
  theory'', \emph{Journal of Sound and Vibration} 258(4), 2002, 649--676.

\bibitem{thomas-shells} O.~Thomas, C.~Touz\'e, A.~Chaigne, ``Non-linear
  vibrations of free-edge thin spherical shells: modal interaction rules
  and 1:1:2 internal resonance'', \emph{International Journal of Solids and
  Structures} 42, 2005, 3339--3373.

\bibitem{nguyen-touze} Q.~Nguyen, C.~Touz\'e, ``Nonlinear vibrations of
  thin plates with variable thickness: application to sound synthesis of
  cymbals'', \emph{Journal of the Acoustical Society of America} 145, 2019,
  977--988.

\bibitem{tolonen} T.~Tolonen, V.~V\"alim\"aki, M.~Karjalainen, ``Modeling
  of tension modulation nonlinearity in plucked strings'', \emph{IEEE
  Transactions on Speech and Audio Processing} 8(3), 2000, 300--310.

\bibitem{bilbao-vk} S.~Bilbao, ``A family of conservative finite
  difference schemes for the dynamical von K\'arm\'an plate equations'',
  \emph{Numerical Methods for Partial Differential Equations} 24(1),
  2008, 193--216.

\bibitem{touze-turbulence} C.~Touz\'e, S.~Bilbao, O.~Cadot, ``Transition
  scenario to turbulence in thin vibrating plates'', \emph{Journal of
  Sound and Vibration} 331(2), 2012, 412--433.

\bibitem{berry} M.~V. Berry, ``Regular and irregular semiclassical
  wavefunctions'', \emph{Journal of Physics A: Mathematical and General}
  10(12), 1977, 2083--2091.

\bibitem{lyon} R.~H. Lyon, R.~G. DeJong, \emph{Theory and Application of
  Statistical Energy Analysis}, 2nd ed., Butterworth-Heinemann, 1995.

\bibitem{zames} G.~Zames, ``On the input-output stability of time-varying
  nonlinear feedback systems --- Part I'', \emph{IEEE Transactions on
  Automatic Control} 11(2), 1966, 228--238.

\bibitem{khalil} H.~K. Khalil, \emph{Nonlinear Systems}, 3rd ed.,
  Prentice Hall, 2002.

\end{thebibliography}

\end{document}
